\clearpage
\begin{proofbox}
\textbf{方法一（最低门槛）：平方差公式 $+$ 割线逼近（不使用导数）}

\textbf{看图理解：}割线 $PQ$ 上的点 $Q$ 沿曲线靠近 $P$，
割线 $PQ$ 的方向就逐渐靠近 $P$ 处切线的方向。
\begin{center}
\textbf{用两点坐标相减求割线斜率}\par\medskip
\begin{tikzpicture}[x=2.3cm,y=2.3cm,>=Stealth,font=\small]
\draw[->,gray] (2.8,1.7)--(5.5,1.7) node[right] {$x$};
\draw[->,gray] (2.8,1.7)--(2.8,4.6) node[above] {$y$};
\draw[blue,very thick,domain=3:5.4,samples=90] plot (\x,{sqrt(3*(\x*\x/4-1))});
\draw[orange!85!black,thick] (3,2)--(5.4,4.4);
\draw[teal!75!black,dashed,thick] plot[domain=3:5.35] (\x,{3+(sqrt(63)/2-3)*(\x-4)});
\draw[gray,dotted] (4,3)--(5,3)--(5,{sqrt(63)/2});
\fill (4,3) circle (1pt) node[above left=5pt] {$P(4,3)$};
\fill (5,{sqrt(63)/2}) circle (1pt) node[above left=4pt] {$Q$};
\node[below] at (4.5,3) {$\Delta x$};
\node[right] at (5,3.5) {$\Delta y$};
\node[orange!85!black,anchor=west] at (3.05,1.86) {切线：$k=1$};
\node[teal!75!black,anchor=west] at (3.05,4.5) {割线 $PQ$：$\Delta y/\Delta x$};
\end{tikzpicture}

\end{center}
图取 $a^2=4,b^2=3$，$P=(4,3)$，$Q=(5,\sqrt{63}/2)$；
$\Delta x,\Delta y$ 是两点坐标之差。下列证明适用于一般的 $a,b>0$。

\textbf{第 1 步：}取双曲线上不同的两点 $P(x_0,y_0)$、$Q(x,y)$：
\[
\frac{x^2}{a^2}-\frac{y^2}{b^2}=1,\qquad
\frac{x_0^2}{a^2}-\frac{y_0^2}{b^2}=1.
\]
两式相减，利用平方差公式 $u^2-v^2=(u-v)(u+v)$，得
\[
\frac{(x-x_0)(x+x_0)}{a^2}
=\frac{(y-y_0)(y+y_0)}{b^2}.
\]
\textbf{第 2 步：}若 $y_0\ne0$，令 $Q$ 在 $P$ 附近且 $Q\ne P$。
此时可以写出割线斜率
\[
 k_{PQ}=\frac{y-y_0}{x-x_0}
 =\frac{b^2(x+x_0)}{a^2(y+y_0)}.
\]
$Q$ 越来越靠近 $P$ 时，$x\to x_0,\ y\to y_0$，
所以割线斜率趋向
\[
 k=\frac{b^2(2x_0)}{a^2(2y_0)}
 =\boxed{\frac{b^2x_0}{a^2y_0}}.
\]
这就是切线斜率。\textbf{注意：}这一步是让两点\emph{逐渐靠近}，
不是在割线斜率 $0/0$ 中直接代入 $Q=P$。

\textbf{第 3 步：}切线经过 $P$，由点斜式得
\[
 y-y_0=\frac{b^2x_0}{a^2y_0}(x-x_0)
 \iff
 \frac{x_0x}{a^2}-\frac{y_0y}{b^2}
 =\underbrace{\frac{x_0^2}{a^2}-\frac{y_0^2}{b^2}}_{=1}.
\]
故切线为 $\boxed{\frac{x_0x}{a^2}-\frac{y_0y}{b^2}=1}$。

\textbf{顶点不能漏：}若 $y_0=0$，由双曲线方程知 $x_0=\pm a$。
用第 1 步的相减结果，得到
\[
 \frac{x-x_0}{y}=\frac{a^2y}{b^2(x+x_0)}\longrightarrow0
 \qquad (Q\to P,\ Q\ne P).
\]
即割线越来越接近竖直方向，所以顶点切线为 $x=x_0$；
代入切点式同样得到 $x_0x/a^2=1$，即 $x=x_0$。

\medskip
\textbf{方法二（最快）：导数 $+$ 点斜式}

\textbf{适用信号：}熟悉导数，且切点 $y_0\ne0$ 时，写两步就够。
在 $P$ 附近，根据 $y_0$ 的正负取对应一段曲线：
\[
 y=\pm\frac{b}{a}\sqrt{x^2-a^2}
 \ \Longrightarrow\ 
 y'=\frac{b^2x}{a^2y},\qquad
 k=\frac{b^2x_0}{a^2y_0}.
\]
立即写切线并整理（$P$ 在双曲线上）：
\[
 y-y_0=\frac{b^2x_0}{a^2y_0}(x-x_0)
 \ \Longrightarrow\ 
 \boxed{\frac{x_0x}{a^2}-\frac{y_0y}{b^2}=1}.
\]
\textbf{例外：}$y_0=0$ 时，不能给含 $y_0$ 的分母代入零；
此时 $x_0=\pm a$，顶点处切线为 $x=x_0$，仍与切点式相符。

\medskip
\textbf{方法三（常规）：直线与双曲线联立 $+$ 判别式}

\textbf{第 1 步：}当 $y_0\ne0$ 时，设过 $P$ 的非竖直直线为
$y=kx+t$，其中 $t=y_0-kx_0$。
代入双曲线并乘以 $a^2b^2$：
\[
 (b^2-a^2k^2)x^2-2a^2ktx-a^2(t^2+b^2)=0.
\]
若 $b^2-a^2k^2\ne0$，这才是关于 $x$ 的一元二次方程。
直线在 $P$ 处相切，$x=x_0$ 应是\textbf{二重根}，因此
\[
 \begin{aligned}
 \Delta
 &=(-2a^2kt)^2+4a^2(b^2-a^2k^2)(t^2+b^2)\\
 &=4a^2b^2(t^2+b^2-a^2k^2)=0.
 \end{aligned}
\]
\textbf{第 2 步：}代入 $t=y_0-kx_0$，再用
$x_0^2-a^2=\frac{a^2y_0^2}{b^2}$、
$y_0^2+b^2=\frac{b^2x_0^2}{a^2}$，得到
\[
 \begin{aligned}
 0&=(y_0-kx_0)^2+b^2-a^2k^2\\
  &=\frac{a^2y_0^2}{b^2}k^2-2x_0y_0k+\frac{b^2x_0^2}{a^2}\\
  &=\left(\frac{ay_0}{b}k-\frac{bx_0}{a}\right)^2.
 \end{aligned}
\]
故 $\displaystyle k=\frac{b^2x_0}{a^2y_0}$，
与前两法一致，代入点斜式即得切点式。
此时 $b^2-a^2k^2=-b^4/y_0^2\ne0$，
也核实了第 1 步使用判别式的前提。

\textbf{顶点单独讨论：}若 $y_0=0$、$x_0=\pm a$，
任何非竖直直线均不能使上述判别式为零；
竖直直线 $x=x_0$ 才是切线，仍满足切点式。

\textbf{结论：}三种方法都得到
$\boxed{\frac{x_0x}{a^2}-\frac{y_0y}{b^2}=1}$。
对于 $\frac{y^2}{a^2}-\frac{x^2}{b^2}=1$，交换 $x$、$y$ 即得
$\boxed{\frac{y_0y}{a^2}-\frac{x_0x}{b^2}=1}$。
\end{proofbox}
